% META: {"id":"1SPE-GEOREP-EX-029","chapitre":"1SPE-GEOMETRIE-REPEREE","type_objet":"exercice","capacites":["1SPE-GEOMETRIE-REPEREE-C3"],"capacites_codes":["C3"],"methodes":["M3"],"parcours":3,"competences":["calculer","raisonner"],"duree_min":15,"mode_creation":"ex_nihilo","sources_inspiration":[],"fichier_tex":"chapitres/1SPE-GEOMETRIE-REPEREE/exercices/1SPE-GEOREP-EX-029.tex","corrige_tex":"chapitres/1SPE-GEOMETRIE-REPEREE/corriges/1SPE-GEOREP-CO-029.tex","status":"approved","origin":"generated","status_history":["generated","audited","accepted","approved"],"release_acceptance":"RELEASE_OWNER_BATCH_ACCEPTANCE","acceptance_closure_digest":"sha256:9b3ccf9a81c5520fbb7e03b7b2d3e7bf2057a02834b6d908bf3be5f49f3b553f","acceptance_receipt_digest":"sha256:4fad86f0282e0fa4e0419cca1ad6379ae7af5a280c04a4a90880f90a1c044242"}
% BEGIN-VERIFY
% from sympy import *
% x, y = symbols('x y')
% # Cercle de centre (a,0) passant par A(1,2) et B(3,-2)
% # (1-a)^2+4 = (3-a)^2+4 => (1-a)^2 = (3-a)^2
% # 1-2a+a^2 = 9-6a+a^2 => 4a = 8 => a = 2
% assert 4*2 == 8
% # r^2 = (1-2)^2 + 4 = 5
% assert (1-2)**2 + 4 == 5
% END-VERIFY
\begin{exercice}{1SPE-GEOREP-EX-029}{3}{15}
Déterminer la valeur de $a$ pour que le cercle de centre $\Omega(a ; 0)$ passe par $A(1 ; 2)$ et $B(3 ; -2)$.
\end{exercice}
